% from ega4-2-fr.tex
\oldpage[IV-2]{10}
\item[\emph{c)}] Le foncteur
$\mathscr G\mapsto\mathscr F\otimes_Y\mathscr G$ de la catégorie des
$\mathcal O_Y$-Modules quasi-cohérents dans celle des
$\mathcal O_X$-Modules quasi-cohérents, est exact et fidèle.
\end{enumerate}
\end{proposition}

Si \emph{b$'$)} est vérifié, il y a au moins un $x\in f^{-1}(y)$ tel que
$(\mathscr F_y)_x\neq0$~; soit $U=\operatorname{Spec}(B)$ un voisinage ouvert
affine de $x$, et soit $\mathscr F|U=\widetilde M$, où $M$ est un $B$-module.
Alors \emph{b$'$)} implique que $M/\mathfrak j_yM$ est $\neq0$, et par suite
(puisque $M$ est un $A$-module plat (2.1.2)) que
$M\otimes_A A_y$ est un $A_y$-module fidèlement plat
$(0_{\mathrm I},6.4.5)$. La relation
$(\mathscr F\otimes_Y\mathscr G)\otimes_Y\mathcal O_y=0$ pour un point fermé
$y$ de $Y$ implique
\[
(\mathscr F\otimes_Y\mathcal O_y)\otimes_{\mathcal O_y}\mathscr G_y=0,
\]
donc $\mathscr G_y=0$. Mais si l'on a $\mathscr G_y=0$ pour tout point fermé
de $Y$, on en conclut $\mathscr G=0$, car si $\mathscr G=\widetilde N$,
l'annulateur d'un élément de $N$ n'est contenu dans aucun idéal maximal de
$A$, donc est $A$ tout entier. La relation
$\mathscr F\otimes_Y\mathscr G=0$ implique donc $\mathscr G=0$, autrement dit
le foncteur $\mathscr G\mapsto\mathscr F\otimes_Y\mathscr G$ est fidèle~; on
sait par ailleurs que ce foncteur est exact (2.1.3), ce qui montre bien que
\emph{b$'$)} entraîne \emph{c)}.

Enfin, pour voir que \emph{c)} entraîne \emph{a)}, on peut se borner au cas
où $Y'$ est aussi affine~; comme $g:Y'\to Y$ est alors un morphisme affine,
il en est de même de la projection $g':X'\to X$ (II, 1.5.5)~; de plus, le
foncteur $\mathscr H'\mapsto g'_*(\mathscr H')$ est alors exact dans la
catégorie des $\mathcal O_{X'}$-Modules quasi-cohérents (I, 1.6.4), et si
$g'_*(\mathscr H')=\mathscr H$, on a $\mathscr H'=\widetilde{\mathscr H}$,
donc le foncteur précédent est aussi fidèle~; pour voir que
$\mathscr G'\mapsto\mathscr F\otimes_Y\mathscr G'$ est exact et fidèle, il
suffit par suite de voir que le foncteur
\[
\mathscr G'\mapsto g'_*(\mathscr F\otimes_Y\mathscr G')
\]
l'est. Or, si $f'=f_{(Y')}:X'\to Y'$, on a
$g'_*(\mathscr F\otimes_Y\mathscr G')
=g'_*(g'^*(\mathscr F)\otimes_{X'}f'^*(\mathscr G'))$~; le fait que $g$ est
affine entraîne que l'on a un isomorphisme canonique
\begin{equation}
\tag{2.2.2.1}
\label{IV.2.2.2.1-fr}
\mathscr F\otimes_{\mathcal O_X}f^*(g_*(\mathscr G'))
\xrightarrow{\sim}
g'_*(g'^*(\mathscr F)\otimes_{\mathcal O_{X'}}f'^*(\mathscr G')).
\end{equation}
En effet, on sait (II, 1.5.2) qu'on a un isomorphisme canonique
\[
f^*(g_*(\mathscr G'))\xrightarrow{\sim}g'_*(f'^*(\mathscr G')),
\]
et d'autre part on a un homomorphisme canonique
$\mathscr F\to g'_*(g'^*(\mathscr F))$ $(0_{\mathrm I},4.4.3.2)$~;
composant l'homomorphisme
$\mathscr F\otimes_{\mathcal O_X}f^*(g_*(\mathscr G'))
\to g'_*(g'^*(\mathscr F))\otimes_{\mathcal O_X}g'_*(f'^*(\mathscr G'))$
avec l'homomorphisme canonique $(0_{\mathrm I},4.2.2.1)$, on en déduit
l'homomorphisme (2.2.2.1)~; la vérification du fait qu'il s'agit d'un
isomorphisme est immédiate en se ramenant au cas où $X$ est affine. Cela
étant, le foncteur $\mathscr G'\mapsto g_*(\mathscr G')$ est exact et fidèle
et par hypothèse il en est de même de
$\mathscr G\mapsto\mathscr F\otimes_Y\mathscr G
=\mathscr F\otimes_{\mathcal O_X}f^*(\mathscr G)$~; leur composé est donc
exact et fidèle, ce qui achève de prouver (2.2.1) et (2.2.2).

\begin{corollary}[2.2.3]
\label{IV.2.2.3-fr}
Soient $X=\operatorname{Spec}(B)$, $Y=\operatorname{Spec}(A)$ deux schémas
affines, $f:X\to Y$ un morphisme, $\mathscr F=\widetilde M$ un
$\mathcal O_X$-Module quasi-cohérent. Pour que $\mathscr F$ vérifie les
conditions équivalentes de (2.2.1) (ou (2.2.2)), il faut et il suffit que le
$A$-module $M$ soit fidèlement plat.
\end{corollary}

En effet, la condition \emph{c)} de (2.2.2) signifie alors que le foncteur
$N\mapsto M\otimes_A N$ de la catégorie des $A$-modules dans celle des
$B$-modules est exact et fidèle, et la conclusion résulte de
$(0_{\mathrm I},6.4.1)$.

\begin{definition}[2.2.4]
\label{IV.2.2.4-fr}
Lorsque les conditions équivalentes de (2.2.1) sont satisfaites, on dit que
le $\mathcal O_X$-Module quasi-cohérent $\mathscr F$ est \emph{fidèlement plat
relativement à} $f$ (ou à $Y$).
\end{definition}

